\section{Discussion}
\subsection{From spectral accessibility to an optimizer}
AQNG is the third step in a sequence. The isotropic readout-rank theory showed that state sensitivity can survive while a low-rank readout retains only a small fraction of visible score information \cite{AitHaddouRank}. The spectral theory generalized that result by replacing global isotropy with relative operator alignment $\Tr(TC_Q)$ \cite{AitHaddouSpectral}. AQNG makes this geometry operational by using the accessible tangent structure to precondition learning.

The empirical results support, but do not prove in full generality, the proposition that useful optimization geometry can be smaller than full state-space geometry. SU(2)-Haar is the clearest example: approximately three quarters of QFIM trace is absent from the $\le2$ accessible metric, yet AQNG generalizes better on average across all four benchmarks. Conversely, the $U(1)$ control shows that reproducing the full-QNG direction is not sufficient for task success. This motivates the hierarchy
\begin{equation}
\boxed{\text{state geometry}\;\longrightarrow\;\text{measurement-accessible geometry}\;\longrightarrow\;\text{task-useful optimization geometry}.}
\end{equation}
The first arrow is governed by measurement/readout contraction. The second additionally depends on loss, labels, data encoding, and finite optimization trajectory.

\subsection{Limitations}
The experiments use six-qubit noiseless statevector simulation, small binary datasets, one fixed split per dataset, and 20 optimization steps. They establish a controlled finite-size comparison, not an asymptotic scaling theorem or hardware advantage. AQNG is not guaranteed to outperform full QNG: RyRz-CZ on Digits is a clear counterexample. The timing comparison is also confounded by different metric simulator backends; a matched-backend study should separately count state preparations, circuit evaluations, shots, and classical linear-algebra cost. The present metric uses fixed diagonal readout probes through weight two; general POVMs, rotated bases, adaptive readouts, finite-shot covariance estimation, and noisy channels can alter accessibility. Finally, the $U(1)$ data do not identify the precise source of task failure. Encoding mismatch, conserved-sector expressivity, and readout/label alignment require targeted ablations.

\section{Conclusion}
We introduced the accessible quantum natural gradient, which replaces the full quantum Fisher metric with the metric induced by measurement/readout-accessible tangent geometry. AQNG is motivated by the spectral identity $\mathbb E[I/F_Q]=\Tr(TC_Q)$ and is therefore an algorithmic application of accessible-tangent theory rather than merely a low-cost QFIM approximation.

Across four binary classification datasets, four six-qubit architectures, and 20 paired seeds per cell, AQNG is competitive with full QNG and improves mean accuracy in nine of twelve nonconserving settings. SU(2)-Haar is particularly informative: the accessible metric retains only about one quarter of QFIM trace while AQNG improves mean accuracy on all four tasks, with corrected significant loss gains on Iris and Wine. Digits/SU(2)-CNOT provides a third corrected significant loss gain. The symmetry-constrained $U(1)$ control remains at fixed chance-like accuracy under both optimizers despite strong direction alignment, separating natural-gradient fidelity from task utility.

These results suggest that variational quantum optimization need not reconstruct every locally distinguishable state-space direction. What matters is the geometry that survives the learning interface and is useful to the task. Establishing when this accessible geometry yields provable optimization or sample-complexity advantages is the central open problem raised by AQNG.

\begin{acknowledgments}
The author thanks the open-source scientific Python and quantum-software communities. Computational results and aggregation files are available in the accompanying reproducibility repository \cite{AQNGRepo}.
\end{acknowledgments}

\section*{Data and code availability}
The complete experiment contains 80 dataset--seed jobs and 640 terminal result rows. Aggregate tables, per-seed results, metric diagnostics, and workflow metadata are stored under \texttt{results/paper/paper\_main\_v1/} in Ref.~\cite{AQNGRepo}. The released manifest records the exact workflow configuration and numerical validation tolerance used for reaggregation.
